\section*{The Propers: Notable and Quotable}

Four documented afterlives of wording from this Sunday's appointed Scripture.
Each pairs a phrase from the Gospel---present in both the longer and the
shorter form---with a later use that redirects it. Straight exegesis and
devotional reception are excluded; those belong to the commentary above. Every
later text was read in a public-domain edition or transcription at the locus
given.

\subsection*{``Patiently bear with the wheat and the tares'' becomes a rule of manners}

\emph{Proper and branch:} Matthew 13:30, both Gospel forms.
\emph{Later use:} Joseph Priestley, \work{An Essay on the First Principles of
Government}, 2nd edition (1771), preface.

Priestley---Dissenting minister, and the chemist who would isolate
oxygen---spends the preface attacking the constitution of church
establishments. Then he turns on his own side: whatever is said against the
abuses, ``let us be particularly careful how we give offence to any serious and
well-disposed minds, and patiently bear with the wheat and the tares growing
together till the harvest.'' He goes straight on to a doctrine of providence in
which ``every thing, even the grossest abuses in the civil or ecclesiastical
constitutions of particular states, is subservient to the wise and gracious
designs'' of God.

The turn is double. First, a command about not uprooting persons becomes a rule
about \emph{tone} in controversy---patience with the prejudices of those
educated inside an institution one is trying to dismantle. Second, the parable
diagnoses the mixture as the work of an enemy, whereas Priestley's providence
absorbs even gross abuse into design; the borrowed sentence and the argument it
introduces pull in opposite directions. \emph{Limits:} Priestley does not cite
chapter and verse, and the phrase governs his manner, not his programme; he
continued to press for disestablishment. The 1771 work is public domain and was
read in the Wikisource transcription of the second edition.

\subsection*{The harvest order reversed}

\emph{Proper and branch:} Matthew 13:29--30, both Gospel forms.
\emph{Later use:} \work{Doctrine and Covenants} 86:6--7, a revelation dated
Kirtland, Ohio, 6 December 1832.

The passage rewrites the parable as an interpretation of church history---the
field ``was the world,'' the apostles were the sowers, the tares were sown
``after they have fallen asleep''---and then changes two things in the Lord's own
words. The reason for delay is transferred from the danger to the wheat to the
weakness of the reapers: ``pluck not up the tares while the blade is yet tender
(for verily your faith is weak), lest you destroy the wheat also.'' And the
order of the harvest is inverted: ``let the wheat and the tares grow together
until the harvest is fully ripe; then ye shall first gather out the wheat from
among the tares, and after the gathering of the wheat, behold and lo, the tares
are bound in bundles, and the field remaineth to be burned.'' Matthew has the
weeds gathered first (13:30), a detail Chrysostom thought was included so that
the wheat would not be alarmed.

The turn is that a restraint grounded in the servants' \emph{ignorance} becomes
a restraint grounded in their \emph{immaturity}, and a sorting that begins with
removal becomes one that begins with rescue. \emph{Limits:} this is a
scriptural claim inside another religious body's canon, and naming a documented
reuse is neither an evaluation of it nor an ecumenical judgment. Text read in
the Wikisource transcription; the 1832 revelation and its nineteenth-century
printings are public domain.

\subsection*{``An enemy hath done this'' applied to a sunk boat}

\emph{Proper and branch:} Matthew 13:28, both Gospel forms (King James
wording).
\emph{Later use:} P.~G. Wodehouse, \work{The White Feather} (1907), chapter 12.

Two schoolboys return from tea to find their hired boat gone from its mooring.
They haul on the rope, the nose of the boat breaks the surface, and Dunstable
delivers judgment: ```There are more things in Heaven and Earth---' said
Dunstable, wiping his hands. `If you ask me, I should say an enemy hath done
this. A boat doesn't sink of its own accord.'\,'' The culprit is a rival
schoolboy named Albert.

The joke works only because the borrowed sentence is grave. In Matthew the
householder's answer explains the origin of evil in the field; in Wodehouse it
explains a scuttled dinghy, and it arrives immediately after a half-finished
line of \emph{Hamlet}, so that the schoolboy's whole rhetorical inheritance is
being spent on vandalism. \emph{Limits:} the wording is the King James
version's, not the Lectionary's, and Wodehouse signals no source; the verbal
match is exact and the sentence is distinctive enough to make dependence
plain, but this is allusion, not quotation with attribution. The 1907 novel is
public domain in the United States and was read in the Wikisource
transcription.

\subsection*{From parable to plant name}

\emph{Proper and branch:} the weeds of Matthew 13:25--30, both Gospel forms.
\emph{Later use:} \work{Easton's Bible Dictionary} (1897), s.v.\ ``Tares'';
\work{Chambers's Twentieth Century Dictionary} (1908), s.v.\ ``Darnel.''

Easton: ``Tares. The bearded darnel, mentioned only in Matt.\ 13:25--30. It is
the Lolium temulentum, a species of rye-grass, the seeds of which are a strong
soporific poison. It bears the closest resemblance to wheat till the ear
appears, and only then the difference is discovered.'' Chambers, a general
English dictionary with no religious purpose, defines the plant by the
parable: ``Darnel \dots\ an annual of the rye-grass genus, the tares of
Scripture,'' with an etymology from Old French \incipit{darne}, stupid, ``from
its supposed narcotic properties.''

The turn is a migration of register. A word that entered English through a
parable becomes the settled common name of a species, and the parable's moral
observation---that the counterfeit cannot be told from the wheat until it
fruits---is restated as a field diagnostic. In Chambers the direction of
definition is fully reversed: the plant is explained by Scripture rather than
Scripture by the plant. That the identified species is a soporific is a
coincidence the lexicographers do not exploit, and this guide notes it as an
observation rather than as their claim. \emph{Limits:} the identification of
\incipit{zizania} with \incipit{Lolium temulentum} is the standard lexical
judgment, not a demonstrated botanical fact of first-century Galilee; both
reference works are public domain and were read in Wikisource transcriptions.
